% Part: history
% Chapter: biographies
% Section: ernst-zermelo

\documentclass[../../../include/open-logic-section]{subfiles}

\begin{document}

\olfileid{his}{bio}{zer}

\olsection{إرنست زيرميلو}

\olphoto{zermelo-ernst}{Ernst Zermelo}

وُلد إرنست زيرميلو في برلين بألمانيا في 27 يوليو 1871. وكانت له خمس
أخوات، لكن أسرته عانت اعتلال الصحة، ولم تبلغ سن الرشد إلا ثلاث من
أخواته. وتوفي والداه أيضًا وهو صغير، فتركا إياه وإخوته أيتامًا حين
كان في السابعة عشرة. وكان لزيرميلو اهتمام عميق بالفنون، ولا سيما
الشعر. وعُرف بحدة الذهن وسرعة البديهة والنزعة النقدية. ومن أشهر
إنجازاته الرياضية إدخال مسلّمة الاختيار عام 1904، وصياغته المسلّمية
لنظرية المجموعات عام 1908.

تنوعت اهتمامات زيرميلو في الجامعة، فدرس الفيزياء والرياضيات والفلسفة.
وأتم، بإشراف هيرمان شفارتس، أطروحته~\emph{بحوث في حساب التغيرات} في
جامعة برلين عام 1894. وقرر عام 1897 مواصلة دراسته في جامعة غوتنغن،
حيث تأثر بشدة بأعمال دافيد هيلبرت في أسس الرياضيات. وفي عام 1899 تأهل
لشغل منصب أستاذ، لكنه لم يحصل على منصب إلا بعد أحد عشر عامًا، وربما
كان ذلك بسبب سلوكه الغريب و«عجلته العصبية».

حصل زيرميلو أخيرًا على منصب أستاذ مدفوع الأجر في جامعة زيورخ عام
1910، لكنه اضطر إلى التقاعد عام 1916 بسبب السل. وبعد تعافيه مُنح
أستاذية فخرية في جامعة فرايبورغ عام 1921. وعمل خلال تلك الفترة على
أسس الرياضيات. وقد أثارت أعمال تورالف سكولم وكورت غودل ضيقه، فانتقد
مناهجهما علنًا في أوراقه. وفُصل من منصبه في فرايبورغ عام 1935 بسبب
افتقاره إلى الشعبية ومعارضته صعود هتلر إلى السلطة في ألمانيا.

اتسمت السنوات الأخيرة من حياة زيرميلو بالعزلة. فبعد فصله عام 1935
هجر الرياضيات، وانتقل إلى الريف حيث عاش عيشة متواضعة. وتزوج عام 1944،
وصار يعتمد اعتمادًا كاملًا على زوجته مع تدهور بصره. وفقد زيرميلو بصره
تمامًا بحلول عام 1951. وتوفي في غونترستال بألمانيا في 21 مايو 1953.

\begin{reading}
للاطلاع على سيرة كاملة لزيرميلو، انظر~\citet{Ebbinghaus2015}. ويمكن
قراءة ورقتيه المؤسستين المنشورتين عامي 1904 و1908 بالألمانية الأصلية
\citep{Zermelo1904,Zermelo1908}. كما تتوفر أعمال زيرميلو المجموعة،
ومنها كتاباته في الفيزياء، بترجمة إنجليزية في
\citep{Ebbinghaus2010,Ebbinghaus2013}.
\end{reading}

\end{document}
